
\documentclass[a4paper,
12pt,footsepline,oneside,			
plainfootsepline,listof=totoc,headinclude,
footinclude]{scrreprt}
\usepackage{graphicx}
\usepackage[utf8]{inputenc}
\usepackage[ngerman]{babel}
\usepackage{geometry}
\usepackage{setspace}
\usepackage[nohyperlinks,printonlyused, withpage,  withpage]{acronym}
\usepackage{float}
\usepackage{tabularx}
\usepackage{wrapfig}
\usepackage{color}
\usepackage{selinput}
\usepackage{fancyvrb}
\usepackage{booktabs}
\usepackage{listings}
\usepackage{blindtext}
\usepackage{scrpage2}
\usepackage{caption}
\usepackage[section]{placeins}
 

\pagestyle{scrheadings}
\usepackage[backend=biber,style=apa]{biblatex}
\DeclareLanguageMapping{german}{german-apa}
\addbibresource{lit.bib}
\captionsetup{justification=raggedright,singlelinecheck=false}

%Seitenzahlen
\clearscrplain
\clearscrheadings
\ofoot[\pagemark]{\pagemark}

%Coding Environment
\definecolor{codegreen}{rgb}{0,0.6,0}
\definecolor{codegray}{rgb}{0.5,0.5,0.5}
\definecolor{codepurple}{rgb}{0.58,0,0.82}
\definecolor{backcolour}{rgb}{0.95,0.95,0.92}
\lstdefinestyle{mystyle}{
	basicstyle= \small\footnotesize\ttfamily,
    backgroundcolor=\color{backcolour},   
    commentstyle=\color{codegreen},
    keywordstyle=\color{magenta},
    numberstyle=\tiny\color{codegray},
    stringstyle=\color{codepurple},
    basicstyle=\footnotesize,
    breakatwhitespace=false,         
    breaklines=true,                 
    captionpos=b,                    
    keepspaces=true,                 
    numbers=left,                    
    numbersep=5pt,                  
    showspaces=false,                
    showstringspaces=false,
    showtabs=false,                  
    tabsize=2
}
\lstset{style=mystyle}

%Umbenennungen
\renewcommand*{\figurename}{Abbildung}
\renewcommand{\contentsname}{Inhaltsverzeichnis}
\renewcommand*{\bibname}{Literaturverzeichnis}
\renewcommand*{\listfigurename}{Abbildungsverzeichnis}

%Absatzformatierung
\renewcommand*\chapterheadstartvskip{\vspace*{-\topskip}}
\renewcommand*\chapterheadendvskip{%
  \vspace*{1\baselineskip plus .1\baselineskip minus .167\baselineskip}}
\parskip 5pt plus 1pt minus 1pt
%\addtolength{\parskip}{\baselineskip}
\parindent 0pt    
  

%Pageoffset
\geometry{
 left=3.5cm,
  right=2.5cm,
  top=2.2cm,
  bottom=3cm,
  bindingoffset=2mm
}

%line spacing
\setstretch{1.5}	


%custom commands
\newcommand{\ThesisTitle}{Text2Process - der Stanford Parser}

\newcommand{\Heading}[1]{ \begin{center} 
\textbf{#1} 
\newline
\end{center}}
 

%document structure
\begin{document}

\pagenumbering{Roman}
\begin{titlepage}
	\centering
%	\includegraphics[width=0.25\textwidth]{pictures/SAP_Logo.png}
	
	\includegraphics{pictures/dhbw_logo.png}
	\vspace{1cm}
	\par
	{\scshape\LARGE Fakultät Wirtschaft\par}
	\vspace{0.5cm}
	{\scshape\Large Studiengang Wirtschaftsinformatik\par}
	\vspace{1cm}
	{\large\bfseries \ThesisTitle \par}
	\vspace{1cm}
	{\Large Seminararbeit "'Integrationsseminar"'\par}

	\vfill
	
	\begin{center}
	\begin{tabularx}{\columnwidth}{XXl}
	Verfasserin: &  \textsc{Jana Kuntz}\\
	Kurs: & \textsc{WWI15B1} \\
	Matrikelnummer: & \textsc{42} \\\\
	Verfasser: & \textsc{Oliver Weisenburger} \\
	Kurs: & \textsc{WWI15B1} \\
	Matrikelnummer: & \textsc{6317874} \\\\
	Seminarbetreuer: &  \textsc{Prof. Dr. Thomas Freytag} 	\\
	Seminargruppe:&  \textsc{Business Processes und Artificial Intelligence} 	 \\
	Abgabedatum: & \textsc{15.01.2018} \\
\end{tabularx} 
 \end{center}

\end{titlepage}
\newpage\cleardoublepage
\Heading{Eidesstattliche Erklärung}
Wir versicheren hiermit, dass wir unsere Seminararbeit mit dem Thema:    \begin{quote}
"\ThesisTitle "
\end{quote} selbstständig verfasst und keine anderen als die angegebenen Quellen und Hilfsmittel benutzt haben. Wir versicheren zudem, dass die eingereichte elektronische Fassung mit der
gedruckten Fassung übereinstimmt.
Diese Arbeit wurde bisher in gleicher oder ähnlicher Form oder auszugsweise
noch keiner Prüfungsbehörde vorgelegt und auch nicht veröffentlicht.

\vspace{50pt} 
\noindent\rule{5cm}{.4pt}\hfill\rule{5cm}{.4pt}\par 
\noindent Datum, Ort  \hspace{7,4cm} Datum, Ort 
\par
\vspace{2cm}
\par
\noindent\rule{5cm}{.4pt}\hfill\rule{5cm}{.4pt}\par 
\noindent Jana Kuntz \hspace{7,4cm} Oliver Weisenburger
\newpage\cleardoublepage

\tableofcontents \newpage\cleardoublepage
\listoffigures  \newpage\cleardoublepage
\addcontentsline{toc}{chapter}{Abkürzungsverzeichnis}
\listoftables  \newpage\cleardoublepage

\chapter*{Abkürzungsverzeichnis}
\begin{acronym} 
  \acro{BPMN}{Business Process Modeling Notation}
\end{acronym}


\chapter{Einleitung}\pagenumbering{arabic}	


\section{example}
Ein Zitat: Der Stanford Parser ist total einfach! (\cite[vgl.][274]{STANFORDNLP}).\par 
Eine Abkürzung: \ac{BPMN}.\par

\begin{figure}[H]
\noindent\includegraphics[width=\linewidth,height=\textheight,
keepaspectratio]{pictures/dhbw_logo.png}
\caption{ein tolles Logo}
\label{fig:DHBW}

Die Abbildung \ref{fig:DHBW} zeigt das DHBW Logo.
\end{figure}

\begin{table}[h!]
  \centering
  \begin{tabular}{cccc}
    \toprule
    Ziele & Gewichtung & Option 1 & Option 2\\
    \midrule
	funktionaler Umfang  	& 40\% & 4 & 5 \\    
	Performance 	& 25\% & 3 & 5 \\    
    Aufwand 		& 15\% & 5 & 2 \\
    Wartbarkeit 	& 10\% & 3 & 3 \\
    Erweiterbarkeit & 10\% & 5 & 3 \\
    \midrule
    Ergebnis 		& 	   & 3,9 & 4,15 \\
    \bottomrule
  \end{tabular}
  \caption{Nutzwertanalyse}
  \label{tab:table1}
\end{table}
\chapter{Hauptteil}


\chapter{Schluss}

\addcontentsline{toc}{chapter}{Literaturverzeichnis}
\printbibliography


\appendix     
\part*{Anhang} % Start the appendix part
   
    
\chapter{Anhang} 

\section{POS-Tags}

\begin{longtable}{|l|l|l|}
\caption{Liste aller POS-Tags der Penn Treebank (\cite[vgl.][8]{PENNTREEBANK})}\\
\hline \multicolumn{1}{|c|}{\textbf{ }} & \multicolumn{1}{c|}{\textbf{Tag}} & \multicolumn{1}{c|}{\textbf{Bedeutung}} \\ \hline 
\endfirsthead

\multicolumn{3}{c}%
{{\bfseries \tablename\ \thetable{} -- Weiterführung}} \\
\hline \multicolumn{1}{|c|}{\textbf{ }} & \multicolumn{1}{c|}{\textbf{Tag}} & \multicolumn{1}{c|}{\textbf{Bedeutung}} \\ \hline 
\endhead

\hline \multicolumn{3}{|r|}{{Weiterführung auf nächster Seite}} \\ \hline
\endfoot

\hline \hline
\endlastfoot
 1. & CC & Coordinating conjunction \\
	 2. & CD & Cardinal number\\
	 3. & DT & Determiner \\
	4.&  EX &	Existential there\\
	5.&	FW &	Foreign word\\
	6.&	IN &	Preposition or subordinating conjunction\\
	7.&	JJ &	Adjective\\
	8.&	JJR &	Adjective, comparative\\
	9.&	JJS &	Adjective, superlative\\
	10.& LS &	List item marker\\
	11.& MD &	Modal\\
	12.& NN &	Noun, singular or mass\\
	13.& NNS &	Noun, plural\\
	14.& NNP &	Proper noun, singular\\
	15.& NNPS &	Proper noun, plural\\
	16.& PDT & Predeterminer\\
	17.& POS & Possessive ending\\
	18. & PRP & Personal pronoun \\
	19.& PRP\$ & Possessive pronoun\\
	20.& RB & Adverb\\
	21.& RBR & Adverb, comparative\\
	22.& RBS & Adverb, superlative\\
	23.& RP & Particle\\
	24.& SYM & Symbol\\
	25.& TO & to\\
	26.& UH & Interjection\\
	27.& VB & Verb, base form\\
	28.& VBD & Verb, past tense\\
	29.& VBG & Verb, gerund or present participle\\
	30.& VBN & Verb, past participle\\
	31.& VBP & Verb, non-3rd person singular present\\
	32.& VBZ & Verb, 3rd person singular present\\
	33.& WDT & Wh-determiner\\
	34.& WP & Wh-pronoun\\
	35.& WP\$ & Possessive wh-pronoun\\
	36.& WRB & Wh-adverb	\\

	  \label{table:POSTAGS}
\end{longtable}



\section{Stanford Dependency-Tags}

\begin{longtable}{|l|l|l|}

  \caption{Liste aller DEP-Tags des Stanford Parsers (\cite[vgl.][3 ff.]{STANFORDDEPS})}\\
\hline \multicolumn{1}{|c|}{\textbf{     }} & \multicolumn{1}{c|}{\textbf{Tag}} & \multicolumn{1}{c|}{\textbf{Bedeutung}} \\ \hline 
\endfirsthead

\multicolumn{3}{c}%
{{\bfseries \tablename\ \thetable{} -- Weiterführung}} \\
\hline \multicolumn{1}{|c|}{\textbf{     }} & \multicolumn{1}{c|}{\textbf{Tag}} & \multicolumn{1}{c|}{\textbf{Bedeutung}} \\ \hline 
\endhead

\hline \multicolumn{3}{|r|}{{Weiterführung auf nächster Seite}} \\ \hline
\endfoot

\hline \hline
\endlastfoot
  1. & acomp & adjectival complement \\
2. & advcl & adverbial clause modifier\\
3. & advmod & adverb modifier\\
4. & agent & agent\\
5. & amod & adjectival modifier\\
6. & appos & appositional modifier\\
7. & aux & auxiliary\\
8. & auxpass & passive auxiliary\\
9. & cc & coordination\\
10. & ccomp & clausal complement\\
11. & conj & conjunct\\
12. & cop & copula\\
13. & csubj & clausal subject\\
14. & csubjpass & clausal passive subject\\
15. & dep & dependent\\
16. & det & determiner\\
17. & discourse & discourse element\\
18. & dobj & direct object\\
19. & expl & expletive\\
20. & goeswith & goes with\\
21. & iobj & indirect object\\
22. & mark & marker\\
23. & mwe & multi-word expression\\
24. & neg & negation modifier\\
25. & nn & noun compound modifier\\
26. & npadvmod & noun phrase as adverbial modifier\\
27. & nsubj & nominal subject\\
28. & nsubjpass & passive nominal subject\\
29. & num & numeric modifier\\
30. & number & element of compound number\\
31. & parataxis & parataxis\\
32. & pcomp & prepositional complement\\
33. & pobj & object of a preposition\\
34. & poss & possession modifier\\
35. & possessive & possessive modifier\\
36. & preconj & preconjunct\\
37. & predet & predeterminer\\
38. & prep & prepositional modifier\\
39. & prepc & prepositional clausal modifier\\
40. & prt & phrasal verb particle\\
41. & punct & punctuation\\
42. & quantmod & quantifier phrase modifier\\
43. & rcmod & relative clause modifier\\
44. & ref & referent\\
45. & root & root\\
46. & tmod & temporal modifier\\
47. & vmod & reduced non-finite verbal modifier\\
48. & xcomp & open clausal complement\\
49. & xsubj & controlling subject\\

\end{longtable}

\section{Beigabenverzeichnis}
\begin{enumerate}
	\item Seminararbeit
	\begin{enumerate}
		\item Seminararbeit Text2WorldModel.pdf
	\end{enumerate}

\end{enumerate} 



\end{document}
